% 02-harmonics.tex -- spherical harmonics.  Formulas F1-F3.
% Input by mostlyharmless.tex.  No preamble here.

\section{Spherical harmonics}
\label{sec:harmonics}

Complex harmonics are labelled by non-negative orders only, as set out in
section~\ref{sec:conventions}, and negative orders are never formed, so no magnitude is ever
needed and there is no signed-order convention to get wrong.  This section records the
harmonic fact the visibility derivation rests on --- orthonormality --- and collects the
real-basis conventions, which the derivation does not use, in one place for comparison with
published work.  Each part says which is argument and which is algebra delegated to a
certificate.

\subsection{Orthonormality}
\label{sec:orthonormality}

With negative orders defined by \eqref{eq:ylm-extension}, the complex harmonics
are orthonormal over the full range,
\begin{equation}
  \int_{S^2} Y_{lm} \, \overline{Y_{l'm'}} \, d\Omega
  = \delta_{ll'} \, \delta_{mm'} , \qquad -l \le m, m' \le l ,
  \label{eq:ylm-orthonormal}
\end{equation}

The proof of \eqref{eq:ylm-orthonormal} is two factors multiplied.  Writing
$d\Omega = \sin\theta\,d\theta\,d\varphi$, the azimuthal factor is the elementary
Fourier statement
\begin{equation}
  \int_0^{2\pi} e^{i(m-m')\varphi} \, d\varphi = 2\pi \, \delta_{mm'} ,
  \label{eq:fourier-orthogonality}
\end{equation}
and the polar factor is
\begin{equation}
  \int_{-1}^{1} P_{lm}(x) \, P_{l'}^m(x) \, dx
  = \delta_{ll'} \, \frac{2}{2l+1} \frac{(l+m)!}{(l-m)!} , \qquad m \ge 0 .
  \label{eq:legendre-orthogonality}
\end{equation}
Equation \eqref{eq:legendre-orthogonality} is a polynomial identity.  Substituting
the Rodrigues form $P_{lm}(x) = (1-x^2)^{m/2} Q_{lm}(x)$ with
$Q_l^m = (-1)^m d^m P_l / dx^m$ turns it into an integral of a polynomial, which
certificate \vtag{1} evaluates exactly for every $0 \le m \le l \le 5$.
\emph{Exact}, range-limited to $l \le 5$.  The negative orders need no separate
treatment: by \eqref{eq:ylm-extension} they are conjugates of the positive ones up
to a sign, and conjugation commutes with the integral.

Multiplying \eqref{eq:fourier-orthogonality} by
\eqref{eq:legendre-orthogonality} gives \eqref{eq:ylm-orthonormal}: the $2\pi$ and
the normalisation together account for the surviving $2/(2l+1)$.  That
multiplication is the argument; the two identities are the certified part.

The implementation is checked end to end on a Gauss--Legendre grid in
$\cos\theta$ and the uniform Fourier rule in $\varphi$.  Both are exact for this
integrand, so the residual is floating point alone.  \emph{Numerical}, $l \le 4$.

The complex harmonics are also compared against SymPy's own \code{Ynm} at $30$
digits.  That is a convention check rather than an identity check, and it is done
numerically because the two expressions are algebraically equal but not
identically written --- a symbolic zero there would test the simplifier more than
the claim.

\subsection{The real basis, and other conventions}
\label{sec:other-conventions}

Nothing in the derivation needs the material in this subsection.  It is collected here so
that the main line stays complex, and so that the conventions most often met in published
work can be translated in one place.

\paragraph{The real basis.}  At degree $l$ there are $2l+1$ real functions, labelled by
$m = 0, \dots, l$ together with a cosine or sine branch:
\begin{equation}
  R_{l,0} = Y_{l0} , \qquad
  R_{l,m,c} = \sqrt{2}\,\mathrm{Re}\,Y_{lm} , \qquad
  R_{l,m,s} = \sqrt{2}\,\mathrm{Im}\,Y_{lm} \quad (m > 0) ,
  \label{eq:real-basis}
\end{equation}
orthonormal and complete,
\begin{equation}
  \int_{S^2} R_{l,m,\sigma} \, R_{l',m',\sigma'} \, d\Omega
  = \delta_{ll'} \, \delta_{mm'} \, \delta_{\sigma\sigma'} ,
  \label{eq:real-orthonormal}
\end{equation}
where $\sigma$ runs over $c$ and $s$, the $m = 0$ function counting once on its own.  A real
map expands in them with real coefficients, $c_{lm}$ and $s_{lm}$, which relate to the stored
half-complex $a_{lm}$ of \eqref{eq:map-from-half-complex} by
\begin{equation}
  a_{l0} = c_{l0} , \qquad
  a_{lm} = \frac{c_{lm} - i\,s_{lm}}{\sqrt{2}} \quad (m > 0) .
  \label{eq:half-complex}
\end{equation}
That $\sqrt{2}$ is the only place the two normalisations differ, and it is the reason a
coefficient quoted from a real-harmonic code has to be rescaled before it means the same
map.

\paragraph{The branches in the visibility.}  The transform of a branch follows from the one
fact about conjugation used throughout, that conjugating a map conjugates its transform and
mirrors the baseline:
\begin{equation}
  V(R_{l,m,c}) \;=\; \frac{V_{lm} + (-1)^m\,\overline{V_{lm}}}{\sqrt2} ,
  \qquad
  V(R_{l,m,s}) \;=\; \frac{V_{lm} - (-1)^m\,\overline{V_{lm}}}{i\sqrt2} ,
  \label{eq:real-branch-visibilities}
\end{equation}
which reduce to the $\cos(m\psi)$ and $\sin(m\psi)$ forms met in the literature.  No object
at a negative order appears in either line.

\paragraph{Signed orders.}  Most of the literature indexes real harmonics by a \emph{signed}
order, with $\cos(m\varphi)$ for $m > 0$ and $\sin(m\varphi)$ for $m < 0$.  That is a
relabelling of \eqref{eq:real-basis}: our $(l, m, c)$ and $(l, m, s)$ are their $(l, m)$ and
$(l, -m)$.  Interferometric surface-mapping codes, including the one the analytic visibility
work we build on uses, follow that convention.  Nothing in this work needs it, but the
translation is what makes published coefficients directly comparable.

\paragraph{The $\sqrt{2}$-free convention.}  Complex coefficients for a real map are
sometimes defined without the $\sqrt{2}$ of \eqref{eq:half-complex}, in which case
\eqref{eq:map-from-half-complex} carries a different numerical factor.  The relation above is
the one used here, and it is the one pyshtools documents for the coefficients of a real
function.

\subsection{What was checked}
\label{sec:harmonics-register}

\begin{center}
\begin{tabular}{l p{0.52\textwidth} l}
\hline
tag & claim & rung \\
\hline
\vtag{1} & Legendre orthogonality & exact, $l \le 5$ \\
         & $\int_{S^2}|Y_{lm}|^2\,d\Omega = 4\pi$ for the implementation & numerical, $l \le 4$ \\
         & the $\sqrt{4\pi}$ against SymPy's orthonormal \code{Ynm} & numerical, $30$ digits \\
\vtag{2} & the extension \eqref{eq:ylm-extension}, and the implementation reproduces it & exact \\
\vtag{3} & the $l = 2$ zonal corollary \eqref{eq:zonal-l2}, against the general \code{sph_harm_addition} & numerical, $40$ digits, $l \le 6$ \\
\vtag{25} & the real basis \eqref{eq:real-basis}: mean square one on the grid, and equal to the complex one & numerical, $l \le 4$ \\
\hline
\end{tabular}
\end{center}
